% =====================================================================
\chapter{Soluciones de los ejercicios}
\label{ap:soluciones}
% =====================================================================

Cada solución se presenta en cinco partes: datos, planteamiento, desarrollo, resultado y comprobación. En los pasos intermedios se redondea a cuatro cifras significativas salvo que se indique otra cosa. Los ejercicios de diseño admiten varias respuestas correctas; se da una razonada.

% ---------------------------------------------------------------------
\section{Tema 1}

\begin{solucion}{eje:1-variables}
\datos reactor con camisa; se miden $T$ y $T_j$; se manipula la válvula; varía la temperatura del agua de red; no se mide $C_A$.
\planteamiento aplicar la \cref{def:ee}: el estado es lo que hay que conocer para predecir; las entradas se dividen en elegidas y no elegidas.
\pasos
\begin{enumerate}[label=\arabic*.]
  \item Estados: $T$, $T_j$ y $C_A$ (cada uno tiene su propio balance, de energía o de materia).
  \item Entrada manipulada: la apertura de la válvula de refrigeración.
  \item Perturbación: la temperatura del agua de red.
  \item Salidas medidas: $T$ y $T_j$.
  \item $C_A$ es un estado no medido: un MPC necesita estimarlo (observador, filtro de Kalman extendido o estimación de horizonte móvil), por ejemplo a partir del balance de energía, porque el calor generado delata la velocidad de reacción.
\end{enumerate}
\resultado \textbf{estados $T,T_j,C_A$; entrada, la válvula; perturbación, el agua de red; salidas, $T$ y $T_j$}. El estado no medido obliga a usar un estimador.
\comprobacion el número de estados coincide con el número de ecuaciones diferenciales del modelo (\cref{def:reactor} más el balance de la camisa).
\end{solucion}

\begin{solucion}{eje:1-zoh}
\datos $R=5$ K/kW; $C=3$ kWh/K; $\tau=15$ h; $\Dt=1$ h.
\planteamiento \cref{prop:zoh} con $A=-1/\tau$ y $B=1/C$ para la entrada $Q$; Euler: $a_E=1-\Dt/\tau$, $b_E=\Dt/C$.
\pasos
\begin{enumerate}[label=\arabic*.]
  \item $a=e^{-1/15}=0{,}93551$.
  \item $b=A^{-1}(a-1)B=(-15)(0{,}93551-1)/3=R(1-a)=5\cdot0{,}06449=0{,}32247$ K/kWh.
  \item $a_E=1-1/15=0{,}93333$; $b_E=1/3=0{,}33333$ K/kWh.
  \item Errores relativos: $(0{,}93333-0{,}93551)/0{,}93551=-0{,}23$~\%; $(0{,}33333-0{,}32247)/0{,}32247=+3{,}37$~\%.
\end{enumerate}
\resultado $\mathbf{a=0{,}9355}$, $\mathbf{b=0{,}3225}$ K/kWh; Euler, $0{,}9333$ y $0{,}3333$, con errores del $-0{,}23$~\% y $+3{,}4$~\%.
\comprobacion la ganancia estática discreta $b/(1-a)=0{,}32247/0{,}06449=5{,}000=R$ K/kW coincide con la continua; con Euler, $b_E/(1-a_E)=0{,}3333/0{,}0667=5$ también: Euler conserva la ganancia estática pero no la dinámica.
\end{solucion}

\begin{solucion}{eje:1-rigidez}
\datos $C_a=0{,}5$, $C_m=5$ kWh/K; $R_{am}=0{,}5$, $R_{ao}=5$ K/kW.
\planteamiento balances de energía de cada nodo; autovalores de una matriz $2\times2$ por traza y determinante; \cref{cor:euler}.
\pasos
\begin{enumerate}[label=\arabic*.]
  \item $C_a\dot T_a=(T_m-T_a)/R_{am}+(\Text-T_a)/R_{ao}$ y $C_m\dot T_m=(T_a-T_m)/R_{am}$. Sin entradas:
  $A=\begin{pmatrix}-(2+0{,}2)/0{,}5&2/0{,}5\\2/5&-2/5\end{pmatrix}=\begin{pmatrix}-4{,}4&4\\0{,}4&-0{,}4\end{pmatrix}$.
  \item $\tr A=-4{,}8$, $\det A=1{,}76-1{,}6=0{,}16$; $\lambda=(-4{,}8\pm\sqrt{23{,}04-0{,}64})/2=(-4{,}8\pm4{,}7329)/2$, es decir, $\lambda_1=-4{,}7664$ y $\lambda_2=-0{,}03357$ h$^{-1}$.
  \item Euler con $\Dt=0{,}25$: $1+\lambda_i\Dt=-0{,}1916$ y $0{,}9916$; ambos de módulo menor que~1: estable.
  \item Euler con $\Dt=0{,}5$: $1-2{,}3832=-1{,}3832$ y $0{,}9832$: inestable.
  \item $\Dt_{\max}=2/4{,}7664=0{,}4196$ h.
\end{enumerate}
\resultado \textbf{$\lambda=-4{,}766$ y $-0{,}0336$ h$^{-1}$; Euler es estable con $\Dt=0{,}25$ h e inestable con $0{,}5$ h; $\Dt_{\max}=0{,}42$ h.}
\comprobacion $\lambda_1+\lambda_2=-4{,}8=\tr A$ y $\lambda_1\lambda_2=0{,}160=\det A$.
\end{solucion}

\begin{solucion}{eje:1-tanque}
\datos $V=2$ m³; $F=0{,}5$ m³/h; $F_1/F=0{,}6$; $UA=1{,}5$ kW/K; $\rho c_p=4180$ kJ/(m³·K); $K_c=2$.
\planteamiento fórmulas del enunciado (lección~5 de 10.450) con unidades coherentes; error permanente de la \cref{prop:offset} con ganancia de lazo $K_3K_c$.
\pasos
\begin{enumerate}[label=\arabic*.]
  \item $\tau_R=V/F=2/0{,}5=4$ h.
  \item $\rho c_pF=4180\cdot0{,}5/3600=0{,}5806$ kW/K; $\beta=1{,}5/0{,}5806=2{,}584$.
  \item $\tau=4/3{,}584=1{,}116$ h; $K_1=0{,}6/3{,}584=0{,}1674$; $K_3=2{,}584/3{,}584=0{,}7210$.
  \item Ganancia de lazo: $K_3K_c=1{,}442$; fracción de error permanente: $1/(1+1{,}442)=0{,}4095$.
\end{enumerate}
\resultado $\mathbf{\tau=1{,}12}$ h, $\mathbf{K_1=0{,}167}$, $\mathbf{K_3=0{,}721}$; con control P queda un $\mathbf{41}$~\% del escalón como error.
\comprobacion $\tau<\tau_R$, como debe: el intercambio de calor acelera la respuesta; y $K_3<1$ porque el vapor solo actúa a través del intercambiador.
\end{solucion}

\begin{solucion}{eje:1-offset}
\datos $k=14$ K/kW; $\Delta r=2$ K.
\planteamiento \cref{prop:offset} y \cref{prop:integral}.
\pasos
\begin{enumerate}[label=\arabic*.]
  \item (a) $e_\infty=2/(1+14\cdot0{,}5)=2/8=0{,}25$ K.
  \item (b) $2/(1+14K_c)<0{,}1\iff1+14K_c>20\iff K_c>1{,}357$ kW/K.
  \item (c) Una ganancia alta satura el actuador ante cualquier cambio, amplifica el ruido de medida y reduce el margen de estabilidad cuando hay retardos; aun así, el error nunca es nulo. La acción integral anula el error en equilibrio para cualquier ganancia que estabilice el lazo, lo que permite elegir $K_c$ por robustez y no por precisión.
\end{enumerate}
\resultado \textbf{(a) $0{,}25$ K; (b) $K_c>1{,}36$ kW/K}; (c) la solución es añadir acción integral, no subir la ganancia.
\comprobacion con $K_c=1{,}357$, $2/(1+19)=0{,}1$ K, en el límite.
\end{solucion}

\begin{solucion}{eje:1-simc}
\datos $k=2$ K/\%; $\tau=10$ min; $\theta=2$ min; $\Dt=0{,}5$ min.
\planteamiento reglas \eqref{eq:simc} y PI posicional de la \cref{def:pid}.
\pasos
\begin{enumerate}[label=\arabic*.]
  \item (a) $\tau_c=2$: $K_c=10/(2\cdot4)=1{,}25$ \%/K; $\tau_I=\min\{10,16\}=10$ min.
  \item (b) $\tau_c=6$: $K_c=10/(2\cdot8)=0{,}625$ \%/K; $\tau_I=\min\{10,32\}=10$ min.
  \item (c) $u_0=K_ce_0+I_0=1{,}25\cdot1+0=1{,}25$ \%; $I_1=I_0+K_c\Dt e_0/\tau_I=1{,}25\cdot0{,}5/10=0{,}0625$ \%.
\end{enumerate}
\resultado \textbf{(a) $K_c=1{,}25$ \%/K, $\tau_I=10$ min; (b) $K_c=0{,}625$ \%/K, $\tau_I=10$ min; (c) $u_0=1{,}25$ \%, $I_1=0{,}0625$ \%.}
\comprobacion duplicar $\tau_c+\theta$ divide $K_c$ por dos, como indica \eqref{eq:simc}; $\tau_I$ queda limitado por $\tau$ en ambos casos.
\end{solucion}

\begin{solucion}{eje:1-integral}
\datos $x_{k+1}=ax_k+bu_k+w$; $u_k=K_c(r-x_k)+I_k$; $I_{k+1}=I_k+\kappa(r-x_k)$ con $\kappa=K_c\Dt/\tau_I$.
\planteamiento escribir el lazo como sistema lineal en $(x,I)$ y estudiar su equilibrio.
\pasos
\begin{enumerate}[label=\arabic*.]
  \item (a) $\begin{pmatrix}x_{k+1}\\I_{k+1}\end{pmatrix}=\begin{pmatrix}a-bK_c&b\\-\kappa&1\end{pmatrix}\begin{pmatrix}x_k\\I_k\end{pmatrix}+\begin{pmatrix}bK_cr+w\\\kappa r\end{pmatrix}$.
  \item (b) Si la matriz tiene radio espectral menor que~1 (\cref{teo:estab}), existe un único equilibrio y el lazo converge a él. En el equilibrio, la segunda ecuación da $I=I+\kappa(r-x)$, luego $x=r$, sea cual sea $w$. La primera fija entonces $I^{*}=\bigl((1-a)r-w\bigr)/b$: el integrador «aprende» la perturbación.
  \item (c) Ejemplo: $a=0{,}9$, $b=1$, $u_{\max}=1$, $w=-0{,}2$, $r=10$. Con $u=1$, el equilibrio de la planta es $x=(1-0{,}2)/(1-0{,}9)=8<10$. El error tiende a $2$ y, sin integración condicional, $I_{k+1}=I_k+2\kappa$ crece sin límite, aunque $u$ siga en~1. Si después la perturbación desaparece, $I$ es enorme y la salida permanece saturada mucho tiempo: sobreoscilación por \emph{windup}.
\end{enumerate}
\resultado \textbf{con lazo estable y sin saturación, el equilibrio cumple $y=r$ para todo $w$; con saturación, el integrador diverge si no hay antisaturación.}
\comprobacion es la \cref{prop:integral} en un caso concreto; en el cuaderno, la comprobación del ejercicio~3 verifica que el integrador no crece con la salida saturada.
\end{solucion}

\begin{solucion}{eje:1-estab}
\datos $\dot x=-x/\tau$; Euler explícito e implícito; modelo rígido del \cref{ej:rigido}.
\planteamiento estudiar el factor de amplificación por paso; para el implícito, diagonalizar.
\pasos
\begin{enumerate}[label=\arabic*.]
  \item (a) Euler: $x_{k+1}=(1-\Dt/\tau)x_k$; $\abs{1-\Dt/\tau}<1\iff0<\Dt/\tau<2\iff\Dt<2\tau$. Si $\tau<\Dt<2\tau$, el factor está en $(-1,0)$ y el signo alterna en cada paso, mientras que la solución exacta $x_0e^{-t/\tau}$ nunca cambia de signo.
  \item (b) Implícito: $(I-\Dt A)x_{k+1}=x_k$. Con $A=V\Lambda V^{-1}$, $(I-\Dt A)^{-1}=V(I-\Dt\Lambda)^{-1}V^{-1}$, de autovalores $1/(1-\lambda_i\Dt)$; si $\lambda_i<0$, $1-\lambda_i\Dt>1$ y el autovalor está en $(0,1)$ para todo $\Dt>0$.
  \item (c) Para un modelo lineal, la discretización exacta ZOH (se calcula $e^{A\Dt}$ una vez); si el modelo es no lineal y rígido, un integrador implícito o la colocación ortogonal, que CasADi y do-mpc ofrecen. Euler explícito obligaría a pasos menores que $0{,}42$ h, dictados por la dinámica rápida del aire, aunque interese la lenta.
\end{enumerate}
\resultado \textbf{Euler explícito: estable si y solo si $\Dt<2\tau$; Euler implícito: estable para todo $\Dt$; para el MPC, ZOH exacto o integradores implícitos.}
\comprobacion con $\tau=0{,}21$ h (dinámica rápida del \cref{ej:rigido}), $2\tau=0{,}42$ h, el mismo límite obtenido con los autovalores.
\end{solucion}

% ---------------------------------------------------------------------
\section{Tema 2}

\begin{solucion}{eje:2-unpaso}
\planteamiento \cref{def:validacion} y \cref{prop:acumula}.
\pasos
\begin{enumerate}[label=\arabic*.]
  \item A un paso, el modelo parte de la medida real en cada instante: sus errores no se acumulan y, en procesos lentos, predecir $T_{k+1}\approx T_k$ ya da un error pequeño.
  \item El MPC encadena 96 predicciones: el error de cada una se suma, amplificado por la dinámica lenta (factor $1/(1-L)$, cercano a 60 en la zona).
  \item Métrica: RMSE (y error máximo) en simulación libre sobre horizontes de la longitud del MPC, en días de validación separados en el tiempo de los de estimación, y con condiciones de operación representativas.
\end{enumerate}
\resultado \textbf{el error a un paso no mide lo que el MPC usa; hay que validar en simulación libre sobre bloques temporales no vistos.}
\comprobacion en el \cref{ej:caso-id}, los tres modelos tienen $0{,}137$ K a un paso y difieren en un factor de casi~2 en simulación.
\end{solucion}

\begin{solucion}{eje:2-arx}
\datos $y=(1;\ 1{,}2;\ 0{,}95;\ 1{,}18)$; $u=(1,0,1)$.
\planteamiento \cref{prop:ls} con $\varphi_k=(y_k,u_k)$ e $Y=(y_1,y_2,y_3)$.
\pasos
\begin{enumerate}[label=\arabic*.]
  \item Filas de $\Phi$: $(1,1)$, $(1{,}2;\ 0)$, $(0{,}95;\ 1)$.
  \item Productos:
  \[
  \Phi\T\Phi=\begin{pmatrix}1+1{,}44+0{,}9025&1+0+0{,}95\\1{,}95&2\end{pmatrix}=\begin{pmatrix}3{,}3425&1{,}95\\1{,}95&2\end{pmatrix},
  \]
  \[
  \Phi\T Y=\begin{pmatrix}1{,}2+1{,}14+1{,}121\\1{,}2+0+1{,}18\end{pmatrix}=\begin{pmatrix}3{,}461\\2{,}38\end{pmatrix}.
  \]
  \item $\det=6{,}685-3{,}8025=2{,}8825$; $\hat a=(2\cdot3{,}461-1{,}95\cdot2{,}38)/2{,}8825=2{,}281/2{,}8825=0{,}7913$; $\hat b=(3{,}3425\cdot2{,}38-1{,}95\cdot3{,}461)/2{,}8825=1{,}2062/2{,}8825=0{,}4185$.
  \item Residuos: $1{,}2-(0{,}7913+0{,}4185)=-0{,}0098$; $0{,}95-0{,}7913\cdot1{,}2=0{,}0004$; $1{,}18-(0{,}7913\cdot0{,}95+0{,}4185)=0{,}0098$.
  \item Simulación desde $y_0=1$: $1{,}2098$; $0{,}7913\cdot1{,}2098=0{,}9573$; $0{,}7913\cdot0{,}9573+0{,}4185=1{,}1760$.
\end{enumerate}
\resultado $\mathbf{\hat a=0{,}791}$, $\mathbf{\hat b=0{,}418}$; simulación $(1;\ 1{,}210;\ 0{,}957;\ 1{,}176)$ frente a los datos $(1;\ 1{,}2;\ 0{,}95;\ 1{,}18)$.
\comprobacion los residuos son ortogonales a las columnas de $\Phi$: $1\cdot(-0{,}0098)+1{,}2\cdot0{,}0004+0{,}95\cdot0{,}0098\approx0$ y $-0{,}0098+0+0{,}0098=0$.
\end{solucion}

\begin{solucion}{eje:2-escalar}
\datos $y_k=\theta_0u_k+e_k$; $\Var e_k=1$; $N=400$.
\planteamiento \cref{prop:ls} con un solo parámetro.
\pasos
\begin{enumerate}[label=\arabic*.]
  \item (a) $\hat\theta_N=\sum u_ky_k/\sum u_k^{2}=\theta_0+\sum u_ke_k/\sum u_k^{2}$. Si $\frac1N\sum u_k^{2}\to\bar\sigma_u^{2}>0$ y $\frac1N\sum u_ke_k\to0$ (independencia), $\hat\theta_N\to\theta_0$.
  \item (b) Condicionado a $u$, $\Var\hat\theta_N=\sum u_k^{2}\Var e/(\sum u_k^{2})^{2}=\lambda/\sum u_k^{2}$.
  \item (c) $u_k=\pm1$: $\sum u_k^{2}=400$ y la desviación típica es $1/\sqrt{400}=0{,}05$. Sinusoide de amplitud $0{,}1$: valor cuadrático medio $0{,}01/2$, $\sum u_k^{2}\approx2$ y desviación típica $1/\sqrt2=0{,}71$.
\end{enumerate}
\resultado \textbf{$\hat\theta_N\to\theta_0$, $\Var\hat\theta_N=\lambda/\sum u_k^{2}$; desviación típica $0{,}05$ frente a $0{,}71$.}
\comprobacion la precisión depende de la energía de la excitación frente al ruido: reducir la amplitud a la décima parte (con la mitad de energía por la forma sinusoidal) multiplica la desviación por $\sqrt{200}=14{,}1$, y $0{,}05\cdot14{,}1=0{,}71$.
\end{solucion}

\begin{solucion}{eje:2-hammerstein}
\datos $y_{k+1}=ay_k+b(u_k+cu_k^{2})$; $\hat\theta=(0{,}85;\ 0{,}30;\ -0{,}06)$.
\planteamiento \cref{ej:hammerstein}.
\pasos
\begin{enumerate}[label=\arabic*.]
  \item (a) $y_{k+1}=ay_k+bu_k+(bc)u_k^{2}$: regresores $(y_k,u_k,u_k^{2})$ y parámetros $(a,b,bc)$, lineal en ellos.
  \item (b) $a=0{,}85$; $b=0{,}30$; $c=-0{,}06/0{,}30=-0{,}2$.
  \item (c) Con $f(u)=\alpha_1u+\alpha_2u^{2}$ libre, $(b,\alpha_1,\alpha_2)$ y $(b/\gamma,\gamma\alpha_1,\gamma\alpha_2)$ dan el mismo producto $bf$ para todo $\gamma\ne0$: solo son identificables $b\alpha_1$ y $b\alpha_2$. Fijar $\alpha_1=1$ elimina la ambigüedad.
\end{enumerate}
\resultado $\mathbf{a=0{,}85}$, $\mathbf{b=0{,}30}$, $\mathbf{c=-0{,}2}$; sin normalización, la identificación es única salvo un factor de escala.
\comprobacion $b\cdot c=0{,}30\cdot(-0{,}2)=-0{,}06$, el tercer parámetro estimado.
\end{solucion}

\begin{solucion}{eje:2-cota}
\datos $L=0{,}98$; $\varepsilon=0{,}01$ K.
\planteamiento \cref{prop:acumula}: $d_k\le\varepsilon(1-L^{k})/(1-L)$.
\pasos
\begin{enumerate}[label=\arabic*.]
  \item $k=4$: $0{,}01(1-0{,}9224)/0{,}02=0{,}0388$ K. $k=24$: $0{,}01(1-0{,}6158)/0{,}02=0{,}1921$ K. $k=96$: $0{,}01(1-0{,}1438)/0{,}02=0{,}4281$ K. Límite: $0{,}5$ K.
  \item (b) $\varepsilon\le0{,}2\cdot0{,}02/(1-0{,}98^{96})=0{,}004/0{,}8562=0{,}00467$ K.
  \item (c) Margen dependiente del paso, $m_j=\varepsilon(1-L^{j})/(1-L)$, que vale cero al principio y crece con el horizonte; o, más sencillo y conservador, un margen constante igual a la cota en el último paso.
\end{enumerate}
\resultado \textbf{$0{,}039$; $0{,}192$; $0{,}428$ y $0{,}5$ K; $\varepsilon\le0{,}0047$ K por paso.}
\comprobacion la cota a un paso es $\varepsilon=0{,}01$ y crece monótonamente hacia $\varepsilon/(1-L)=0{,}5$.
\end{solucion}

\begin{solucion}{eje:2-lazo}
\datos $a=0{,}9$; $b=0{,}5$; $K=0{,}4$.
\planteamiento \cref{prop:lazo}.
\pasos
\begin{enumerate}[label=\arabic*.]
  \item (a) $y_{k+1}=(0{,}9-0{,}5\cdot0{,}4)y_k+\varepsilon_{k+1}=0{,}7y_k+\varepsilon_{k+1}$.
  \item (b) Cualquier par con $a'-0{,}4b'=0{,}7$: $(0{,}7;\ 0)$, $(0{,}9;\ 0{,}5)$, $(1{,}1;\ 1)$.
  \item (c) Con $u_k=-0{,}4y_k+v_k$: $\frac1N\sum y_k^{2}\to\sigma_y^{2}$, $\frac1N\sum y_ku_k\to-0{,}4\sigma_y^{2}$ (porque $y_k$ solo depende de $v$ hasta $k-1$ y $v_k$ es independiente de $y_k$) y $\frac1N\sum u_k^{2}\to0{,}16\sigma_y^{2}+\sigma_v^{2}$. El determinante del límite es $\sigma_y^{2}(0{,}16\sigma_y^{2}+\sigma_v^{2})-0{,}16\sigma_y^{4}=\sigma_y^{2}\sigma_v^{2}>0$.
\end{enumerate}
\resultado \textbf{se observa $y_{k+1}=0{,}7y_k$; los datos no distinguen $(a,b)$; con excitación externa la matriz límite es invertible.}
\comprobacion el par verdadero $(0{,}9;\ 0{,}5)$ está entre los del apartado (b), y el determinante se anula si $\sigma_v=0$, que es el caso de la \cref{prop:lazo}.
\end{solucion}

\begin{solucion}{eje:2-sesgo}
\datos $x_{k+1}=ax_k+bu_k$; $y_k=x_k+v_k$; $u$, $v$ blancos e independientes.
\planteamiento límite del estimador de mínimos cuadrados: $\theta^{*}=R_{\varphi\varphi}^{-1}r_{\varphi y}$.
\pasos
\begin{enumerate}[label=\arabic*.]
  \item Varianza estacionaria: $\sigma_x^{2}=a^{2}\sigma_x^{2}+b^{2}\sigma_u^{2}$, luego $\sigma_x^{2}=b^{2}\sigma_u^{2}/(1-a^{2})$.
  \item $\E[y_k^{2}]=\sigma_x^{2}+\sigma_v^{2}$; $\E[y_ku_k]=0$, porque $x_k$ depende de $u$ hasta $k-1$; $\E[u_k^{2}]=\sigma_u^{2}$. La matriz $R_{\varphi\varphi}$ es diagonal.
  \item $\E[y_ky_{k+1}]=\E[(x_k+v_k)(ax_k+bu_k+v_{k+1})]=a\sigma_x^{2}$; $\E[u_ky_{k+1}]=b\sigma_u^{2}$.
  \item Por tanto $\hat\alpha\to a\sigma_x^{2}/(\sigma_x^{2}+\sigma_v^{2})$ y $\hat\beta\to b$.
  \item Con $a=0{,}9$, $b=\sigma_u=\sigma_v=1$: $\sigma_x^{2}=1/0{,}19=5{,}263$; $\hat\alpha\to0{,}9\cdot5{,}263/6{,}263=0{,}7563$.
  \item Constantes de tiempo, en muestras: la verdadera es $-1/\ln0{,}9=9{,}49$ y la estimada, $-1/\ln0{,}7563=3{,}58$.
\end{enumerate}
\resultado \textbf{$\hat\alpha\to a\,\sigma_x^{2}/(\sigma_x^{2}+\sigma_v^{2})$: atenuación hacia cero; el modelo es casi tres veces más rápido que la planta.}
\comprobacion si $\sigma_v=0$, $\hat\alpha\to a$; el ruido en un regresor atenúa su coeficiente, el fenómeno de error en las variables, y es la razón de preferir la validación en simulación libre y los métodos de error de salida.
\end{solucion}

\begin{solucion}{eje:2-diseno}
\planteamiento aplicar el tema completo: excitación (\cref{def:pe}), lazo cerrado (\cref{prop:lazo}), partición (\cref{obs:particion}), validación (\cref{def:validacion}), gemelo (\cref{def:gemelo}).
\pasos
\begin{enumerate}[label=\arabic*.]
  \item Históricos del PID: útiles para la parte estacional y las perturbaciones, pero poco informativos sobre la respuesta a la entrada (\cref{prop:lazo}). Se usan como prueba final por estación, no para estimar.
  \item Experimentos: dos semanas con escalones aleatorios de consigna o de potencia, de 1 a 4~h, dentro de una banda segura (por ejemplo, $\pm2$ K respecto a la consigna habitual y solo fuera del horario de ocupación), con parada automática si se sale de la banda de confort ampliada.
  \item Partición: semana~1 para estimar, semana~2 para validar; los históricos de otra estación, como prueba.
  \item Modelos: ARX y NARX con regresores físicos (producto potencia por temperatura exterior, radiación), modelo RC de dos nodos ajustado por mínimos cuadrados no lineales y, si hay datos, un modelo de estado neuronal.
  \item Métricas: RMSE y error máximo en simulación libre a 24~h, sesgo medio y error en la energía diaria.
  \item Criterio de aceptación: por ejemplo, RMSE en simulación a 24~h menor que $0{,}3$ K y error de energía diaria menor que el 5~\%; el margen del MPC se fija con el error observado (\cref{eje:2-cota}).
  \item Recalibración: vigilar la media diaria del residuo (\cref{eje:6-deriva}) y reestimar con ventana móvil cuando supere el umbral, validando el modelo nuevo en sombra antes de sustituir al anterior.
\end{enumerate}
\resultado \textbf{experimento acotado y excitante en lazo abierto, partición temporal, validación en simulación libre con umbrales, y recalibración disparada por el seguimiento de residuos.}
\comprobacion cada paso responde a uno de los errores frecuentes de la \cref{obs:errores2}.
\end{solucion}

% ---------------------------------------------------------------------
\section{Tema 3}

\begin{solucion}{eje:3-deslizante}
\planteamiento \cref{def:mpc}.
\pasos
\begin{enumerate}[label=\arabic*.]
  \item El plan es óptimo para la predicción hecha con la información de $k$; en $k+1$ hay una medida nueva que incorpora perturbaciones y errores de modelo. Volver a resolver es lo que convierte el plan en realimentación.
  \item En lazo abierto, un error de $2$ K en la temperatura exterior desplaza la temperatura interior durante horas (\cref{prop:acumula}) y el plan, calculado para otra trayectoria, provocaría violaciones de confort o calentamiento innecesario sin corregirlos.
\end{enumerate}
\resultado \textbf{se aplica solo la primera decisión para incorporar cada nueva medida; en lazo abierto los errores de previsión no se corrigen.}
\comprobacion el MPC del cuaderno mantiene el confort pese al ruido de medida y al error del modelo identificado, precisamente porque reoptimiza en cada paso.
\end{solucion}

\begin{solucion}{eje:3-lqr}
\datos $A=2$; $B=1$; $Q=R=1$; $P_f=0$; $N=2$; $x_0=1$.
\planteamiento recursión del \cref{teo:lqr}.
\pasos
\begin{enumerate}[label=\arabic*.]
  \item $P_2=0$; $K_1=(1+0)^{-1}\cdot0=0$; $P_1=1+4\cdot0-0=1$.
  \item $K_0=(1+1)^{-1}\cdot1\cdot2=1$; $P_0=1+4\cdot1-2\cdot1\cdot1=3$.
  \item Secuencia: $u_0=-x_0=-1$, $x_1=2-1=1$, $u_1=0$, $x_2=2$. Coste: $1+1+1+0=3=P_0x_0^{2}$.
  \item Horizonte infinito: $P=1+4P-4P^{2}/(1+P)$, es decir, $P^{2}-4P-1=0$, $P=2+\sqrt5=4{,}236$, $K=2P/(1+P)=1{,}618$, y el lazo cerrado tiene autovalor $2-1{,}618=0{,}382$.
\end{enumerate}
\resultado \textbf{$P_1=1$, $K_1=0$, $P_0=3$, $K_0=1$; coste $3$; $P=4{,}236$, $K=1{,}618$, autovalor $0{,}382$.}
\comprobacion el sistema abierto es inestable ($\abs{2}>1$) y el LQR de horizonte infinito lo estabiliza ($0{,}382<1$); con horizonte~2 y $P_f=0$, el último paso no actúa ($K_1=0$) porque el estado final no cuesta nada.
\end{solucion}

\begin{solucion}{eje:3-condensada}
\datos $a=0{,}9$; $b=0{,}5$; $N=3$; $Q=P_f=1$; $R=0{,}1$.
\planteamiento \eqref{eq:condensada} y \cref{prop:qp}.
\pasos
\begin{enumerate}[label=\arabic*.]
  \item $\Phi=(0{,}9;\ 0{,}81;\ 0{,}729)\T$; $\Gamma=\begin{pmatrix}0{,}5&0&0\\0{,}45&0{,}5&0\\0{,}405&0{,}45&0{,}5\end{pmatrix}$.
  \item $\Gamma\T\Gamma$: diagonal $0{,}25+0{,}2025+0{,}1640=0{,}6165$; $0{,}25+0{,}2025=0{,}4525$; $0{,}25$; fuera de la diagonal $0{,}225+0{,}18225=0{,}40725$, $0{,}2025$ y $0{,}225$.
  \item $H=\Gamma\T\Gamma+0{,}1I=\begin{pmatrix}0{,}7165&0{,}4073&0{,}2025\\0{,}4073&0{,}5525&0{,}225\\0{,}2025&0{,}225&0{,}35\end{pmatrix}$.
  \item Autovalores: $0{,}185$; $0{,}273$; $1{,}161$, todos positivos.
\end{enumerate}
\resultado \textbf{$H\succ0$, con autovalores entre $0{,}185$ y $1{,}161$.}
\comprobacion el menor autovalor es mayor que $R=0{,}1$, como garantiza $v\T Hv\ge v\T\bar Rv$.
\end{solucion}

\begin{solucion}{eje:3-bateria}
\datos precios $(0{,}10;\ 0{,}30;\ 0{,}05;\ 0{,}25)$; $\eta=0{,}9$; capacidad y potencia de $1$.
\planteamiento energía que vuelve a la red: $\eta^{2}$ por kWh comprado.
\pasos
\begin{enumerate}[label=\arabic*.]
  \item (a) Comprar $1$ kWh cuesta $p_c$ y devuelve $\eta^{2}$ kWh, vendidos a $p_v$: beneficio $p_v\eta^{2}-p_c>0\iff p_v/p_c>1/\eta^{2}$.
  \item (b) Ciclo~1: $0{,}30\cdot0{,}81-0{,}10=0{,}143$; ciclo~2: $0{,}25\cdot0{,}81-0{,}05=0{,}1525$; total $0{,}2955$ €. Con $(0{,}20;\ 0{,}24)$: $0{,}24/0{,}20=1{,}2<1{,}235$ y $0{,}24\cdot0{,}81-0{,}20=-0{,}0056$: el segundo ciclo no se hace y el beneficio es $0{,}143$ €.
  \item (c) Cada ciclo mueve $1+0{,}81=1{,}81$ kWh por la red y cuesta $0{,}0362$ € de degradación: $0{,}1068+0{,}1163=0{,}2231$ €.
\end{enumerate}
\resultado \textbf{umbral $1/\eta^{2}=1{,}235$; $0{,}2955$ €; $0{,}143$ € con el segundo par cambiado; $0{,}2231$ € con degradación.}
\comprobacion con $\eta=1$ y sin degradación se recupera $0{,}40$ € del \cref{ej:bateria-lp}.
\end{solucion}

\begin{solucion}{eje:3-precalentar}
\datos $\tau=15$ h; 4~h entre las 6 y las 10; precios $0{,}10$ y $0{,}22$ €/kWh.
\planteamiento el exceso de temperatura decae como $e^{-t/\tau}$ (\cref{ej:zoh}); precio efectivo $=$ precio $/$ fracción retenida.
\pasos
\begin{enumerate}[label=\arabic*.]
  \item (a) $e^{-4/15}=0{,}766$.
  \item (b) $0{,}10/0{,}766=0{,}131<0{,}22$ €/kWh: conviene.
  \item (c) Con $\tau=5$ h: $e^{-0{,}8}=0{,}449$ y $0{,}10/0{,}449=0{,}223>0{,}22$: ya no conviene (por poco).
  \item Con COP: el calor cuesta $0{,}10/2{,}4=0{,}0417$ €/kWh$_{\mathrm{t}}$ a las 6 y $0{,}22/2{,}9=0{,}0759$ a las 10. Con $\tau=15$ h el efectivo es $0{,}0417/0{,}766=0{,}0544<0{,}0759$: conviene; con $\tau=5$ h, $0{,}0417/0{,}449=0{,}0927>0{,}0759$: no conviene.
\end{enumerate}
\resultado \textbf{se retiene el 77~\%; precalentar conviene en el edificio pesado y no en el ligero.}
\comprobacion es lo que hace el MPC del \cref{ej:caso-mpc}: almacenar calor en la masa térmica solo compensa si la constante de tiempo es larga frente al desfase entre valle y punta.
\end{solucion}

\begin{solucion}{eje:3-blanda}
\datos $J=3+4u_0+2u_1+3u_0^{2}+2u_0u_1+2u_1^{2}$; $u_0\ge-0{,}5-s$, $s\ge0$; multiplicador de la restricción dura $0{,}5$.
\planteamiento con la holgura activa, la condición de estacionariedad en $u_0$ es $\partial J/\partial u_0=\rho$.
\pasos
\begin{enumerate}[label=\arabic*.]
  \item $\rho=0{,}2$: $4+6u_0+2u_1=0{,}2$ y $2+2u_0+4u_1=0$. De la segunda, $u_1=-(1+u_0)/2$; sustituyendo, $5u_0=-2{,}8$, $u_0=-0{,}56$, $u_1=-0{,}22$, $s=0{,}06$.
  \item Coste: $J=1{,}604$ y $J+\rho s=1{,}616$.
  \item $\rho=1\ge0{,}5$: la holgura no compensa; $s=0$, $u=(-0{,}5;\ -0{,}25)$ y $J=1{,}625$, la solución dura.
\end{enumerate}
\resultado \textbf{con $\rho=0{,}2$ la restricción se viola en $0{,}06$; con $\rho=1$ se recupera exactamente la solución dura.}
\comprobacion \cref{prop:exacta}: el umbral es el multiplicador $0{,}5$; resolviendo con CVXPY se obtienen los mismos valores.
\end{solucion}

\begin{solucion}{eje:3-riccati}
\planteamiento inducción sobre la recursión y comparación de costes de horizontes distintos.
\pasos
\begin{enumerate}[label=\arabic*.]
  \item Simetría: $P_N=P_f$ es simétrica. Si $P_{k+1}$ lo es, $A\T P_{k+1}A$ es simétrica y $A\T P_{k+1}BK_k=A\T P_{k+1}BM^{-1}B\T P_{k+1}A$ también, porque $M=R+B\T P_{k+1}B$ es simétrica; luego $P_k$ lo es.
  \item Semidefinición: $x\T P_kx=V_k(x)$ es el mínimo de una suma de términos $\ge0$ (\cref{teo:lqr}).
  \item Monotonía con $P_f=0$: $V_{k-1}(x)$ es el coste óptimo con un paso más que $V_k(x)$. Dada la secuencia óptima de $V_{k-1}$, sus primeros $N-k$ pasos forman una secuencia admisible para el horizonte menor, con coste no mayor (se suprimen términos no negativos); luego $V_k(x)\le V_{k-1}(x)$.
\end{enumerate}
\resultado \textbf{$P_k=P_k\T\succeq0$ y, con $P_f=0$, $x\T P_kx$ crece al aumentar el horizonte restante.}
\comprobacion en el \cref{eje:3-lqr}, $P_2=0\le P_1=1\le P_0=3$, que convergen hacia $P=4{,}236$ al alargar el horizonte.
\end{solucion}

\begin{solucion}{eje:3-terminal}
\datos $x^{+}=x+u$; $\ell=x^{2}+u^{2}$; $\abs u\le0{,}5$; $P=1{,}618$; $K=0{,}618$.
\planteamiento comprobar las hipótesis del \cref{teo:mpc}.
\pasos
\begin{enumerate}[label=\arabic*.]
  \item (a) Con $u=-Kx$, $x^{+}=(1-K)x=0{,}382x$. $V_f(x^{+})-V_f(x)=1{,}618(0{,}1459-1)x^{2}=-1{,}382x^{2}$ y $\ell=(1+0{,}382)x^{2}=1{,}382x^{2}$: se cumple con igualdad, por ser $P$ solución de Riccati.
  \item (b) $\abs{Kx}\le0{,}5\iff\abs x\le0{,}5/0{,}618=0{,}809$. Si $\abs x\le0{,}809$, $\abs{x^{+}}=0{,}382\abs x\le0{,}809$: invariante. $\Xset_f=[-0{,}809;\ 0{,}809]$.
  \item (c) $f(0,0)=0$, $\ell\ge x^{2}$, y $V_f$, $\kappa_f$ y $\Xset_f$ cumplen las hipótesis; por el \cref{teo:mpc}, para todo $N\ge1$ y todo $x_0$ desde el que $\mathbb P_N$ sea factible (los que pueden llegar a $\Xset_f$ en $N$ pasos, $\abs{x_0}\le0{,}809+0{,}5N$), el lazo cerrado es factible siempre y $x_k\to0$.
\end{enumerate}
\resultado \textbf{la condición de decrecimiento se cumple con igualdad; $\Xset_f=[-0{,}809;\ 0{,}809]$; estabilidad para todo $N$ en la región factible.}
\comprobacion $1{,}618\cdot0{,}382^{2}+1+0{,}618^{2}=0{,}236+1+0{,}382=1{,}618=P$: la ecuación de Riccati escalar.
\end{solucion}

% ---------------------------------------------------------------------
\section{Tema 4}

\begin{solucion}{eje:4-truncado}
\datos $r\approx-0{,}05$ €; $\gamma=0{,}99$; truncación a los 288 pasos.
\planteamiento valor estacionario $r/(1-\gamma)$; objetivo $y=r+\gamma(1-d)V(s')$.
\pasos
\begin{enumerate}[label=\arabic*.]
  \item (a) $V\approx-0{,}05/0{,}01=-5$ €.
  \item (b) Tratada como terminación: $y=-0{,}05$. Correctamente ($d=0$): $y=-0{,}05+0{,}99\cdot(-5)=-5{,}0$.
  \item (c) Los estados cercanos al corte parecen cien veces mejores de lo que son. Si la observación permite reconocerlos (por ejemplo, por la hora), el agente aprende que al final no importa lo que haga y deja de calentar; si no puede reconocerlos, el objetivo es ruidoso e inconsistente.
\end{enumerate}
\resultado \textbf{$V\approx-5$ €; objetivo $-0{,}05$ (incorrecto) frente a $-5{,}0$ (correcto).}
\comprobacion $-0{,}05+0{,}99\cdot(-5)=-5$, consistente con un valor estacionario: el estado siguiente vale lo mismo que el actual.
\end{solucion}

\begin{solucion}{eje:4-constante}
\planteamiento sumar la constante dentro del retorno.
\pasos
\begin{enumerate}[label=\arabic*.]
  \item (a) $G_t'=\sum_k\gamma^{k}(R_{t+k+1}+c)=G_t+c/(1-\gamma)$. Tomando esperanzas, $V'^{\pi}=V^{\pi}+c/(1-\gamma)$ y $Q'^{\pi}=Q^{\pi}+c/(1-\gamma)$ para toda política; el desplazamiento es igual para todas las acciones y no cambia el $\arg\max$.
  \item (b) Desde $s$, «terminar» da recompensa $0$ y acaba; «seguir» da $-0{,}5$ y en el paso siguiente acaba con recompensa~$0$. Con $\gamma=1$: terminar vale $0$ y seguir $-0{,}5$; es mejor terminar. Con $c=1$: terminar vale $1$ y seguir $0{,}5+1=1{,}5$; ahora es mejor seguir.
\end{enumerate}
\resultado \textbf{en tareas continuadas, $V'=V+c/(1-\gamma)$ y la política no cambia; en episódicas, la constante premia o castiga la duración y puede cambiarla.}
\comprobacion en (b) el retorno con $c$ suma $c$ multiplicado por el número de pasos, que depende de la política: es exactamente lo que no ocurre en la tarea continuada.
\end{solucion}

\begin{solucion}{eje:4-sobre}
\datos $Q=(0;\ -0{,}5)$; $\sigma=1$.
\planteamiento $\E[\max(X,Y)]$ para gaussianas independientes: con $\vartheta=\sqrt{\sigma_X^{2}+\sigma_Y^{2}}$ y $\alpha=(\mu_X-\mu_Y)/\vartheta$, vale $\mu_X\Phi(\alpha)+\mu_Y\Phi(-\alpha)+\vartheta\varphi(\alpha)$.
\pasos
\begin{enumerate}[label=\arabic*.]
  \item (a) $\vartheta=\sqrt2=1{,}414$; $\alpha=0{,}354$; $\Phi(-\alpha)=0{,}362$; $\varphi(\alpha)=0{,}375$. $\E[\max]=0-0{,}5\cdot0{,}362+1{,}414\cdot0{,}375=0{,}349>0$.
  \item (b) $\E[\min(\hat Q_1,\hat Q_2)]=-\sigma/\sqrt\pi=-0{,}564$, por simetría con la \cref{prop:sobre}.
  \item (c) Con un solo crítico, los valores $Q$ estimados tienden a superar a los retornos observados y el error crece por el \emph{bootstrapping}; con el mínimo de dos críticos, el sesgo pasa a ser negativo y los valores $Q$ quedan por debajo o cerca de los retornos, lo que se ve al comparar las curvas de $Q$ y de retorno de evaluación en la tarea.
\end{enumerate}
\resultado \textbf{sesgo $+0{,}349$ con el máximo; $-0{,}564$ con el mínimo de dos estimaciones.}
\comprobacion si las dos acciones valieran lo mismo, la fórmula daría $\vartheta\varphi(0)=1{,}414\cdot0{,}399=0{,}564=1/\sqrt\pi$, el caso de la \cref{prop:sobre}.
\end{solucion}

\begin{solucion}{eje:4-blanda}
\datos $Q=(1,2,0)$; $\alpha\in\{1;\ 0{,}1\}$.
\planteamiento \cref{prop:blanda}.
\pasos
\begin{enumerate}[label=\arabic*.]
  \item $\alpha=1$: $Z=e+e^{2}+1=11{,}107$; $\pi^{*}=(0{,}2447;\ 0{,}6652;\ 0{,}0900)$; $V=\ln11{,}107=2{,}4076$.
  \item $\E[Q]=0{,}2447+1{,}3305=1{,}5752$; $\mathcal H=0{,}8324$ nats; suma $2{,}4076$.
  \item $\alpha=0{,}1$: $\pi^{*}=(4{,}5\cdot10^{-5};\ 0{,}99995;\ 2\cdot10^{-9})$; $V=0{,}1\ln(e^{10}+e^{20}+1)=2{,}0000045$; $\mathcal H=0{,}0005$ nats.
\end{enumerate}
\resultado \textbf{$V=2{,}408$ con $\alpha=1$ y $V=2{,}000005$ con $\alpha=0{,}1$.}
\comprobacion al bajar la temperatura, la política se concentra en la acción de mayor $Q$ y $V\to\max Q=2$.
\end{solucion}

\begin{solucion}{eje:4-temperatura}
\datos $\mathcal H=-1{,}5$; $\bar{\mathcal H}=-1$; $\alpha=0{,}2$; tasa $0{,}1$; $\sigma=0{,}5$; $u=0{,}5$.
\planteamiento $\partial J/\partial\alpha=\E[-\log\pi]-\bar{\mathcal H}=\mathcal H-\bar{\mathcal H}$; entropía gaussiana $\frac12\ln(2\pi e\sigma^{2})$; corrección $-\log(1-\tanh^{2}u)$.
\pasos
\begin{enumerate}[label=\arabic*.]
  \item (a) $\partial J/\partial\alpha=-1{,}5+1=-0{,}5$; $\alpha\leftarrow0{,}2-0{,}1\cdot(-0{,}5)=0{,}25$.
  \item (b) $\mathcal H=\frac12\ln(2\pi e\cdot0{,}25)=\frac12\ln4{,}270=0{,}7258$ nats.
  \item $\tanh0{,}5=0{,}4621$; $\ln(1-0{,}4621^{2})=\ln0{,}7864=-0{,}2402$; luego $\log\pi(a)=\log\mathcal N(u)+0{,}2402$.
\end{enumerate}
\resultado \textbf{$\alpha$ sube a $0{,}25$; $\mathcal H=0{,}726$ nats; la compresión suma $0{,}240$ al logaritmo de la densidad.}
\comprobacion la entropía estaba por debajo del objetivo y la temperatura sube para premiar más la exploración; la densidad aumenta tras $\tanh$ porque la función comprime el intervalo.
\end{solucion}

\begin{solucion}{eje:4-lagrange}
\datos $(J_r,J_c)$: $\pi_1=(-18;\ 2)$, $\pi_2=(-19;\ 0{,}5)$, $\pi_3=(-21;\ 0)$; $d=1$; $\eta=0{,}5$.
\planteamiento $L_i(\lambda)=J_r-\lambda(J_c-d)$; ascenso dual $\lambda\leftarrow\max\{0,\lambda+\eta(J_c-d)\}$.
\pasos
\begin{enumerate}[label=\arabic*.]
  \item (a) Factibles: $\pi_2$ y $\pi_3$; la mejor es $\pi_2$ ($-19$). Si se admiten mezclas, $\frac13\pi_1+\frac23\pi_2$ tiene $J_c=1$ y $J_r=-18{,}67$: el óptimo de un CMDP puede ser aleatorizado~\cite{altman}.
  \item (b) $L_1=-18-\lambda$, $L_2=-19+0{,}5\lambda$, $L_3=-21+\lambda$. $L_1\ge L_2\iff\lambda\le2/3$; $L_2\ge L_3\iff\lambda\le4$. Gana $\pi_1$ si $\lambda<2/3$, $\pi_2$ si $2/3<\lambda<4$ y $\pi_3$ si $\lambda>4$.
  \item (c) $\lambda_0=0\to\pi_1$, $\lambda_1=0{,}5\to\pi_1$, $\lambda_2=1\to\pi_2$, $\lambda_3=0{,}75\to\pi_2$, $\lambda_4=0{,}5\to\pi_1$: oscila entre $\pi_1$ y $\pi_2$.
\end{enumerate}
\resultado \textbf{óptima determinista $\pi_2$; el ascenso dual oscila entre $\pi_1$ y $\pi_2$ y, en promedio, se comporta como una mezcla de ambas.}
\comprobacion la mezcla del apartado (a) tiene la restricción activa, $2/3+1/3=1$, y su valor $-18{,}67$ queda entre los de $\pi_1$ y $\pi_2$, como corresponde a un multiplicador en $[2/3,4]$.
\end{solucion}

\begin{solucion}{eje:4-shaping}
\datos $\Phi(s)=-\kappa\max\{\underline T-T,0\}$; $\kappa=1$; $\gamma=0{,}99$; $T=19\to19{,}5$ °C; $\underline T=20$ °C.
\planteamiento repetir la suma telescópica del \cref{teo:shaping} con horizonte aleatorio $T_f$.
\pasos
\begin{enumerate}[label=\arabic*.]
  \item (a) Si el episodio termina en $s_{T_f}$ terminal, $\sum_{t<T_f}\gamma^{t}F_t=\gamma^{T_f}\Phi(s_{T_f})-\Phi(s_0)=-\Phi(s_0)$, porque $\Phi$ es nula en los terminales. Como antes, $Q'=Q-\Phi(s)$ y el $\arg\max$ no cambia.
  \item (b) $\Phi(s)=-1$, $\Phi(s')=-0{,}5$; $F=0{,}99\cdot(-0{,}5)-(-1)=0{,}505>0$: se premia acercarse a la banda.
  \item (c) Una bonificación $+0{,}1$ cada vez que $T$ sube, sin penalización cuando baja: el agente puede calentar y dejar enfriar en ciclo para cobrarla indefinidamente. No deriva de un potencial: su suma sobre un ciclo cerrado no es nula.
\end{enumerate}
\resultado \textbf{la política óptima se conserva con terminales si $\Phi=0$ en ellos; $F=0{,}505$; las bonificaciones no potenciales crean ciclos explotables.}
\comprobacion sobre cualquier ciclo $s_0\to\dots\to s_0$, con $\gamma=1$ la suma de $F$ de un potencial es nula, y la de la bonificación de (c) es positiva.
\end{solucion}

\begin{solucion}{eje:4-simlema}
\datos $r_{\max}=1$; $\epsilon=0{,}01$.
\planteamiento \cref{teo:simlema}; descomposición de la pérdida de rendimiento en tres términos.
\pasos
\begin{enumerate}[label=\arabic*.]
  \item (a) $\gamma=0{,}9$: $0{,}9\cdot0{,}01/0{,}01=0{,}9$, frente a un rango de $V$ de anchura $2/(1-\gamma)=20$ (4,5~\%). $\gamma=0{,}99$: $0{,}99\cdot0{,}01/10^{-4}=99$, frente a 200 (50~\%): casi inútil.
  \item (b) $V^{\pi^{*}}-V^{\hat\pi^{*}}=(V^{\pi^{*}}-\hat V^{\pi^{*}})+(\hat V^{\pi^{*}}-\hat V^{\hat\pi^{*}})+(\hat V^{\hat\pi^{*}}-V^{\hat\pi^{*}})$. El término central es $\le0$ porque $\hat\pi^{*}$ es óptima en el simulador; cada uno de los otros está acotado por el lema. Pérdida $\le2\gamma\epsilon r_{\max}/(1-\gamma)^{2}$.
  \item (c) El error crece como $(1-\gamma)^{-2}$: con un modelo imperfecto conviene un horizonte efectivo más corto o, mejor, replanificar con realimentación frecuente como hace el MPC, que reinicia la predicción desde el estado medido.
\end{enumerate}
\resultado \textbf{cota de $0{,}9$ con $\gamma=0{,}9$ y de $99$ con $\gamma=0{,}99$; pérdida de la política del simulador $\le2\gamma\epsilon r_{\max}/(1-\gamma)^{2}$.}
\comprobacion con $\epsilon=0$ la cota es nula, como debe: modelos idénticos dan valores idénticos.
\end{solucion}

% ---------------------------------------------------------------------
\section{Tema 5}

\begin{solucion}{eje:5-kpi}
\planteamiento protocolo de la sección~5.1 y errores de la \cref{obs:errores5}.
\pasos
\begin{enumerate}[label=\arabic*.]
  \item ¿Respecto a qué controlador y con qué ajuste?
  \item ¿Mismo periodo, o antes-después con tiempo meteorológico y uso distintos? ¿Se normalizó por grados-día y ocupación?
  \item ¿Se mantuvo el confort? ¿Con qué indicador?
  \item ¿Energía o coste? ¿De qué consumo: climatización o edificio entero?
  \item ¿Cuántos días, con qué dispersión e intervalo de confianza?
  \item ¿Tenían ambos controladores la misma información (previsiones, sensores)?
\end{enumerate}
\resultado \textbf{sin referencia, protocolo, confort e incertidumbre, un porcentaje de ahorro no es interpretable.}
\comprobacion cada pregunta corresponde a uno de los cinco requisitos del protocolo.
\end{solucion}

\begin{solucion}{eje:5-ahorro}
\datos \cref{tab:hvac}; $\rho=2$ €/(K·h).
\planteamiento ahorro $=1-\text{valor}/\text{valor del PI}$.
\pasos
\begin{enumerate}[label=\arabic*.]
  \item (a) Energía: MPC $1-171{,}1/175{,}4=2{,}5$~\%; RL $1-169{,}9/175{,}4=3{,}1$~\%. Coste: MPC $1-19{,}28/21{,}06=8{,}5$~\%; RL $1-20{,}63/21{,}06=2{,}0$~\% (la tabla da $8{,}4$ y $2{,}1$~\% porque el cuaderno los calcula con los valores sin redondear).
  \item (b) Total del MPC sin margen: $18{,}76+2\cdot3{,}54=25{,}84$ €, peor que $19{,}28$ €. Sería preferible si $18{,}76+3{,}54\rho<19{,}28$, es decir, $\rho<0{,}147$ €/(K·h).
  \item (c) El MPC desplaza el consumo de punta a valle usando la masa térmica: la energía apenas cambia, el precio medio pagado sí.
\end{enumerate}
\resultado \textbf{MPC: $-2{,}5$~\% de energía y $-8{,}4$~\% de coste; RL: $-3{,}1$~\% y $-2{,}0$~\%; el MPC sin margen solo compensaría con $\rho<0{,}15$ €/(K·h).}
\comprobacion precio medio pagado: PI $21{,}06/175{,}4=0{,}120$ €/kWh; MPC $19{,}28/171{,}1=0{,}113$ €/kWh.
\end{solucion}

\begin{solucion}{eje:5-crn}
\datos $\sigma=2$ €; $r=0{,}9$.
\planteamiento \cref{prop:crn}.
\pasos
\begin{enumerate}[label=\arabic*.]
  \item (a) Independientes: $\sqrt{2\cdot4}=2{,}83$ €; comunes: $\sqrt{2\cdot4\cdot0{,}1}=0{,}894$ €.
  \item (b) Con $n$ pares independientes, $2{,}83/\sqrt n=0{,}894\Rightarrow n=8/0{,}8=10$.
\end{enumerate}
\resultado \textbf{$2{,}83$ € frente a $0{,}89$ €; una semana común equivale a diez independientes.}
\comprobacion el cociente de varianzas es $1/(1-r)=10$.
\end{solucion}

\begin{solucion}{eje:5-semenov}
\datos los del \cref{ej:caso-reactor}; $Q_g(350)=20$ kW.
\planteamiento \cref{prop:semenov}: $Q_g(T)=20\exp\bigl(10^{4}(1/350-1/T)\bigr)$ y $Q_g'=Q_gE/(RT^{2})$.
\pasos
\begin{enumerate}[label=\arabic*.]
  \item (a) $10^{4}(1/350-1/360)=0{,}7937$; $Q_g(360)=20e^{0{,}7937}=44{,}23$ kW; $Q_g'=44{,}23\cdot10^{4}/360^{2}=3{,}413$ kW/K. Haría falta $UA>3{,}41$ kW/K, con $T_j=360-44{,}23/UA$.
  \item (b) $UA=1{,}5$: $T_j=350-20/1{,}5=336{,}67$ K, pero $Q_g'(350)=1{,}633>1{,}5$: el equilibrio es inestable.
  \item (c) Bajar la temperatura de operación hasta donde $Q_g'<1{,}5$ kW/K, que ocurre por debajo de $348{,}9$ K (con $T_j=336{,}7$ K), con margen adicional; dosificar el reactivo (semilotes) para reducir $C_A$; y programar la limpieza del intercambiador.
\end{enumerate}
\resultado \textbf{$Q_g'(360)=3{,}41$ kW/K; con $UA=1{,}5$ kW/K el punto de $350$ K se vuelve inestable.}
\comprobacion la pendiente crece con $T$ como $Q_g/T^{2}$, exponencialmente, mientras que la retirada es lineal: subir $T$ siempre acerca al límite de estabilidad.
\end{solucion}

\begin{solucion}{eje:5-pue}
\datos reducción del 40~\% en refrigeración y del 15~\% en el exceso de PUE; PUE ilustrativo $1{,}12$.
\planteamiento exceso $=$ PUE $-1$; la refrigeración es una fracción $f$ del exceso.
\pasos
\begin{enumerate}[label=\arabic*.]
  \item (a) $0{,}40f=0{,}15\Rightarrow f=0{,}375$.
  \item (b) Exceso $0{,}12\to0{,}12\cdot0{,}85=0{,}102$; PUE $1{,}102$. Consumo total: $(1{,}12-1{,}102)/1{,}12=1{,}6$~\%.
  \item (c) Los tres porcentajes (40, 15 y 1,6) se refieren a bases distintas: refrigeración, exceso sobre la carga informática y consumo total.
\end{enumerate}
\resultado \textbf{la refrigeración era el 37,5~\% del exceso; PUE $1{,}102$; un 1,6~\% del consumo total.}
\comprobacion $0{,}375\cdot0{,}12\cdot0{,}40=0{,}018=1{,}12-1{,}102$.
\end{solucion}

\begin{solucion}{eje:5-bateria}
\datos previsión $(0{,}10;\ 0{,}30;\ 0{,}05;\ 0{,}25)$; precios reales $(0{,}10;\ 0{,}30;\ 0{,}28;\ 0{,}25)$; $1$ kWh, $\eta=1$.
\planteamiento plan del \cref{ej:bateria-lp}: cargar, descargar, cargar, descargar.
\pasos
\begin{enumerate}[label=\arabic*.]
  \item (a) Lazo abierto: $-0{,}10+0{,}30-0{,}28+0{,}25=0{,}17$ €.
  \item (b) MPC: hora~1, carga; hora~2, descarga a $0{,}30$; hora~3, observa $0{,}28$ y, con la previsión $0{,}25$ para la hora~4, comprar perdería $0{,}03$: espera. Beneficio $0{,}20$ €.
  \item (c) Información perfecta: $0{,}20$ € (un solo ciclo; mantener la carga hasta la hora~3 daría $0{,}18$).
\end{enumerate}
\resultado \textbf{$0{,}17$ € en lazo abierto, $0{,}20$ € con el MPC, igual al óptimo con información perfecta.}
\comprobacion la diferencia ($0{,}03$ €) es exactamente la pérdida del segundo ciclo, que el MPC evita al observar el precio real antes de decidir.
\end{solucion}

\begin{solucion}{eje:5-semenov-dem}
\planteamiento condiciones de tangencia; convexidad de $h(T)=Q_g(T)-UA(T-T_j)$.
\pasos
\begin{enumerate}[label=\arabic*.]
  \item (a) Tangencia: $Q_g(T^{*})=UA(T^{*}-T_j)$ y $Q_g'(T^{*})=UA$. Como $Q_g'=Q_gE/(RT^{2})$, la segunda da $Q_g(T^{*})=UA\,RT^{*2}/E$; sustituyendo en la primera, $T^{*}-T_j=RT^{*2}/E$.
  \item (b) $Q_g''=Q_g\bigl[(E/RT^{2})^{2}-2E/(RT^{3})\bigr]>0$ si $T<E/(2R)$ (5000~K aquí): $h$ es convexa. Si $h(T_1)=h(T_2)=0$ con $T_1<T_2$ y no hay más ceros, $h<0$ en $(T_1,T_2)$; luego $h'(T_1)\le0\le h'(T_2)$, y las desigualdades son estrictas porque, si $h'(T_1)=0$, $T_1$ sería un mínimo de la función convexa y $h\ge0$ en todas partes. Por la \cref{prop:semenov}, $T_1$ es estable y $T_2$ inestable.
  \item (c) $Q_g$ es proporcional a $C_A$, que solo disminuye durante el lote: congelarla en su valor inicial sobrestima el calor generado y su pendiente en cada temperatura, de modo que si el criterio se cumple con $C_{A,0}$ también se cumple después.
\end{enumerate}
\resultado \textbf{$T_j^{*}=T^{*}-RT^{*2}/E$; con dos equilibrios, el inferior es estable y el superior inestable; la hipótesis de $C_A$ constante es conservadora.}
\comprobacion con los datos del \cref{ej:caso-reactor}: $T^{*}=352{,}69$ K y $352{,}69-352{,}69^{2}/10^{4}=340{,}26$ K; y $Q_g'(350)=1{,}633<2<2{,}414=Q_g'(355{,}23)$.
\end{solucion}

\begin{solucion}{eje:5-protocolo}
\planteamiento \cref{prop:crn}, protocolo de la sección~5.1 y \cref{def:validar}.
\pasos
\begin{enumerate}[label=\arabic*.]
  \item Asignación: días alternos o bloques de dos días, sorteando el orden dentro de cada par (para no confundir con el ciclo semanal), durante al menos 8 semanas.
  \item Indicadores: energía y coste de climatización diarios, grados-hora fuera de banda en ocupación, potencia máxima y número de intervenciones del operador.
  \item Normalización: por grados-día de calefacción y por ocupación, o mejor, un modelo de regresión del consumo con temperatura exterior, radiación y ocupación como covariables y el controlador como factor.
  \item Análisis: contraste $t$ pareado sobre los pares de días (\cref{ej:contraste}) o el modelo de regresión; intervalo de confianza del ahorro.
  \item Parada: volver al PID si la violación diaria supera un umbral (por ejemplo, $1$ K·h en ocupación), si hay quejas o si el MPC falla en resolver; todo registrado.
\end{enumerate}
\resultado \textbf{diseño emparejado y aleatorizado, indicadores de energía, coste y confort normalizados, contraste pareado e intervalos, y criterios de parada definidos de antemano.}
\comprobacion es el método publicado de activar y desactivar el sistema que se usó en las pruebas del centro de datos~\cite{evansgao}, con un análisis estadístico explícito.
\end{solucion}

% ---------------------------------------------------------------------
\section{Tema 6}

\begin{solucion}{eje:6-modos}
\planteamiento \cref{def:modos}.
\pasos
\begin{enumerate}[label=\arabic*.]
  \item (a) Asesor: el sistema no ejecuta. (b) Supervisado: ejecuta dentro de límites con el operador vigilando. (c) Autónomo con supervisión: actúa cada cinco minutos con salvaguardas automáticas y los operadores pueden salir del modo de IA~\cite{gamblegao}.
  \item Evidencia para pasar de (a) a (b): un periodo en sombra con las recomendaciones registradas; tasa de aceptación por el operador; ninguna recomendación que hubiera violado límites; ahorro estimado con un experimento pareado; filtro y respaldo probados.
\end{enumerate}
\resultado \textbf{(a) asesor; (b) supervisado; (c) autónomo con supervisión.}
\comprobacion la secuencia coincide con la evolución publicada del caso de Google, de recomendaciones (2016) a control autónomo (2018).
\end{solucion}

\begin{solucion}{eje:6-capas}
\planteamiento \cref{fig:capas}: cada función en la capa de su escala de tiempo.
\pasos
\begin{enumerate}[label=\arabic*.]
  \item Temperatura de impulsión: control básico (PID), segundos.
  \item Consignas horarias: control supervisor (MPC o RL), minutos.
  \item Parada por alta presión: protección o sistema de seguridad, milisegundos a segundos, con lógica y sensores propios.
  \item Previsión del precio: planificación, horas o días.
  \item Detección de sonda averiada: supervisión y diagnóstico (comparación con redundancias o con el gemelo), minutos, con alarma al operador.
  \item La protección no puede depender de la capa supervisora porque esta es software complejo, puede fallar o comunicarse mal y, por diseño, no se certifica como función de seguridad; la protección debe actuar aunque todo lo demás falle.
\end{enumerate}
\resultado \textbf{básico, supervisor, seguridad, planificación y diagnóstico, respectivamente; la seguridad es independiente.}
\comprobacion la asignación respeta que cada capa trabaje más lenta que la inferior.
\end{solucion}

\begin{solucion}{eje:6-alarmas}
\datos 9 alarmas/h; referencias de ISA-18.2: unas 150/día (1 cada 10 min) aceptable; unas 300/día (2 cada 10 min) máximo manejable.
\planteamiento convertir tasas.
\pasos
\begin{enumerate}[label=\arabic*.]
  \item (a) $9\cdot24=216$ al día; $9/6=1{,}5$ cada 10 minutos: entre lo aceptable y el máximo manejable.
  \item (b) $12\cdot24=288$ al día; $2$ cada 10 minutos: en el máximo manejable, sin margen para episodios anómalos.
  \item (c) Racionalizar: solo alarmas que exijan acción (la salida al modo de respaldo y la pérdida de comunicación sí; un residuo algo alto no); agrupar las relacionadas; prioridades; suprimir las del MPC cuando está desactivado; enviar lo informativo a registros y paneles.
\end{enumerate}
\resultado \textbf{216 y 288 al día; la capa nueva lleva la carga al máximo manejable si no se racionaliza.}
\comprobacion $288/300=0{,}96$ del máximo de referencia.
\end{solucion}

\begin{solucion}{eje:6-contraste}
\datos $n=20$; $\bar d=1{,}20$ €; $s=1{,}50$ €; $t_{0{,}975;19}=2{,}093$.
\planteamiento contraste $t$ pareado (módulo~0.9, §5.3).
\pasos
\begin{enumerate}[label=\arabic*.]
  \item (a) $s/\sqrt{20}=0{,}3354$; $t=3{,}578$; $p=0{,}0020$; intervalo $[0{,}498;\ 1{,}902]$ €.
  \item (b) No significativo si $t<2{,}093$: $s>1{,}2\sqrt{20}/2{,}093=2{,}564$ €.
  \item (c) Semiamplitud $t_{0{,}975;n-1}\cdot1{,}5/\sqrt n\le0{,}3$: con la aproximación normal, $n\approx(1{,}96\cdot1{,}5/0{,}3)^{2}=96$; con la $t$, $n=99$ días.
  \item (d) El emparejamiento elimina la variabilidad entre días (tiempo meteorológico, ocupación) y evita confundir el efecto del controlador con las diferencias entre inviernos.
\end{enumerate}
\resultado \textbf{$t=3{,}58$, $p=0{,}002$, $[0{,}50;\ 1{,}90]$ €; no significativo si $s>2{,}56$ €; unos 99 días para $\pm0{,}3$ €.}
\comprobacion el intervalo no contiene el cero, coherente con $p<0{,}05$.
\end{solucion}

\begin{solucion}{eje:6-deriva}
\datos $\sigma=0{,}14$ K; $n=96$.
\planteamiento error típico de la media $\sigma/\sqrt n$.
\pasos
\begin{enumerate}[label=\arabic*.]
  \item (a) $0{,}14/\sqrt{96}=0{,}0143$ K; umbral $3\cdot0{,}0143=0{,}043$ K.
  \item (b) Los residuos de un sistema lento están muy autocorrelacionados: el número efectivo de muestras independientes es mucho menor que 96, el error típico real es mayor y el umbral de (a) daría falsas alarmas frecuentes. Hay que estimar la varianza de la media diaria con datos de días sin deriva (medias por lotes) o con la autocorrelación.
  \item (c) Por ejemplo: alarma de recalibración si la media diaria supera el umbral corregido tres días seguidos, o un CUSUM sobre las medias diarias; nunca recalibrar automáticamente sin validar el modelo nuevo en sombra.
\end{enumerate}
\resultado \textbf{error típico ingenuo $0{,}014$ K y umbral $0{,}043$ K, que hay que ampliar por la autocorrelación.}
\comprobacion con autocorrelación positiva $\Var\bar e=\frac{\sigma^{2}}{n}\bigl(1+2\sum_{j\ge1}(1-j/n)r_j\bigr)>\sigma^{2}/n$.
\end{solucion}

\begin{solucion}{eje:6-filtro}
\datos $a=0{,}9835$; $b=0{,}23$; $d\in[-0{,}0827;\ 0{,}4133]$; $[17,26]$ °C; $P\in[0,3]$.
\planteamiento fórmula de $\Uset_{\mathrm{seg}}(T)$ del \cref{ej:filtro}.
\pasos
\begin{enumerate}[label=\arabic*.]
  \item (a) $T=20$: $aT=19{,}67$; inferior $(17-19{,}67+0{,}0827)/0{,}23<0\to0$; superior $(26-19{,}67-0{,}4133)/0{,}23=25{,}7\to3$: $[0,3]$. $T=25{,}9$: $aT=25{,}473$; superior $(26-25{,}473-0{,}4133)/0{,}23=0{,}496$: $[0;\ 0{,}496]$.
  \item (b) A $25{,}9$ °C se aplica $0{,}496$ kW (recorte de $0{,}5$); a $17$ °C, $\Uset_{\mathrm{seg}}=[1{,}579;\ 3]$ y se aplica $1{,}579$ kW.
  \item (c) Con $\Text\le35$ °C, $d_{\max}=0{,}01653\cdot45=0{,}744$; en $T=26$, el extremo superior es $(26-25{,}571-0{,}744)/0{,}23=-1{,}37<0$: $\Uset_{\mathrm{seg}}$ vacío. Ninguna acción garantiza seguir en el intervalo cuando $T>25{,}68$ °C (sin refrigeración); $[17,26]$ deja de ser invariante.
\end{enumerate}
\resultado \textbf{$[0,3]$ y $[0;\ 0{,}496]$; se aplican $0{,}496$ y $1{,}579$ kW; con $35$ °C de máxima el intervalo no es invariante.}
\comprobacion en $T=25{,}9$ con $P=0{,}496$ y $d=d_{\max}$: $25{,}473+0{,}114+0{,}413=26{,}000$ °C, justo en el límite.
\end{solucion}

\begin{solucion}{eje:6-invariancia}
\planteamiento los extremos de $\Uset_{\mathrm{seg}}(x)$ son funciones afines decrecientes de $x$.
\pasos
\begin{enumerate}[label=\arabic*.]
  \item (a) $\Uset_{\mathrm{seg}}(x)=[\max\{0,\ell(x)\},\min\{\bar u,m(x)\}]$ con $\ell(x)=(x_{\min}-ax-w_{\min})/b$ y $m(x)=(x_{\max}-ax-w_{\max})/b$. Es no vacío si y solo si (i) $\ell(x)\le\bar u$, (ii) $m(x)\ge0$ y (iii) $\ell(x)\le m(x)$.
  \item (iii) equivale a $w_{\max}-w_{\min}\le x_{\max}-x_{\min}$, independiente de $x$. (i) es más exigente cuanto menor es $x$ (porque $\ell$ decrece): basta en $x_{\min}$, $(1-a)x_{\min}-w_{\min}\le b\bar u$. (ii) es más exigente cuanto mayor es $x$: basta en $x_{\max}$, $(1-a)x_{\max}\ge w_{\max}$. Luego no vacío en todo el intervalo $\iff$ no vacío en sus extremos, y por la \cref{prop:invariancia} el intervalo es invariante.
  \item (b) Con $\abs{u_k-u_{k-1}}\le\delta$, las acciones admisibles dependen también de $u_{k-1}$: el estado relevante es $(x_k,u_{k-1})$ y la invariancia debe comprobarse en ese espacio aumentado; puede ocurrir que $x$ sea seguro pero, viniendo de una acción muy distinta, no haya tiempo de corregir.
\end{enumerate}
\resultado \textbf{invariante si y solo si $(1-a)x_{\min}-w_{\min}\le b\bar u$, $(1-a)x_{\max}\ge w_{\max}$ y $w_{\max}-w_{\min}\le x_{\max}-x_{\min}$.}
\comprobacion con los datos del \cref{ej:filtro}: $0{,}0165\cdot17+0{,}0827=0{,}363\le0{,}69$; $0{,}0165\cdot26=0{,}429\ge0{,}413$; $0{,}496\le9$.
\end{solucion}

\begin{solucion}{eje:6-plan}
\planteamiento recorrer las secciones del tema y enlazar con los anteriores.
\pasos
\begin{enumerate}[label=\arabic*.]
  \item Arquitectura: el MPC escribe consignas horarias a los PID de la bomba de calor por OPC UA, a través del filtro; el SIS y las protecciones de la máquina no cambian (sección~6.1).
  \item Modos: dos semanas en sombra, cuatro en asesor, después supervisado; criterios de avance cuantitativos (sección~6.2, \cref{eje:6-modos}).
  \item Alarmas nuevas: salida a respaldo, pérdida de comunicación y fallo del resolvedor, racionalizadas (sección~6.3).
  \item Filtro: intervalo de la \cref{prop:invariancia}, límite de velocidad, perro guardián con transferencia sin saltos al PI (\cref{prop:bumpless}).
  \item Validación: experimento pareado de 8 semanas y contraste (\cref{ej:contraste}, \cref{eje:5-protocolo}).
  \item Seguimiento: residuos diarios del gemelo con umbral corregido y recalibración validada en sombra (\cref{eje:6-deriva}, \cref{tema:2}).
  \item Ciberseguridad: zona propia para la capa supervisora, conducto autenticado y registrado, operación sin nube (sección~6.6).
  \item Regulación: análisis de si el sistema es componente de seguridad de una infraestructura crítica (anexo~III del Reglamento de IA); en una oficina, normalmente no; documentar la supervisión humana igualmente (sección~6.7).
\end{enumerate}
\resultado \textbf{un plan por etapas con salvaguardas, validación estadística, seguimiento y análisis de ciberseguridad y regulación.}
\comprobacion cada elemento corresponde a un error frecuente de la \cref{obs:errores6}.
\end{solucion}
